% Generated draft: requires mathematical and source review.
\documentclass[11pt]{article}
\usepackage[margin=25mm]{geometry}
\usepackage{fontspec}
\setmainfont{Latin Modern Roman}
\usepackage{amsmath,amssymb,mathtools,unicode-math}
\setmathfont{Latin Modern Math}
\usepackage{enumitem,needspace,xurl,hyperref}
\setlength{\parindent}{0pt}
\setlength{\parskip}{5pt}
\setlength{\emergencystretch}{3em}
\newcommand{\N}{\mathbb{N}}
\newcommand{\Z}{\mathbb{Z}}
\newcommand{\Q}{\mathbb{Q}}
\newcommand{\R}{\mathbb{R}}
\newcommand{\C}{\mathbb{C}}
\newcommand{\F}{\mathbb{F}}
\newcommand{\A}{\mathbb{A}}
\newcommand{\PP}{\mathbb{P}}
\newcommand{\E}{\mathbb{E}}
\newcommand{\eps}{\varepsilon}
\newcommand{\dd}{\,\mathrm{d}}
\begin{document}
\section*{UnsolvedMath ID: 30001618}
\textbf{Source problem code:} catalogue entry
\subsection{Definitions and mathematical statement}
\textbf{Terminology.} A one-dimensional interacting particle system describes particles whose stochastic evolution depends on their configuration in one spatial dimension. In an exclusion system, each site can contain at most one particle. Finite-range exclusion permits jumps only over a bounded spatial distance; nearest-neighbour exclusion restricts jumps to adjacent sites. The source does not specify the underlying spatial domain, jump rates, or allowed interactions.

Integrated particle-density fluctuations refer to a spatially accumulated, centered particle-density field, with a reference point and a time-dependent current term where required by the chosen height convention. The centering, reference convention, and initial distribution are not specified in the supplied record. KPZ denotes the Kardar--Parisi--Zhang universality class, and EW denotes the Edwards--Wilkinson universality class. The record does not specify a mathematical definition of membership in either class.

\textbf{Source research problem.} Determine explicit hypotheses on one-dimensional interacting particle systems extending beyond nearest-neighbour simple exclusion, including finite-range exclusion systems, under which their integrated particle-density fluctuations rigorously exhibit the predicted KPZ or EW scaling behaviour. The desired results should specify which hypotheses lead to each class and establish the corresponding fluctuation limits.

\textbf{Unresolved formulation.} For a precise convergence assertion, one must specify the stochastic dynamics, spatial domain and boundary conditions, initial laws or stationary measures, density regime, centering and observation frame, space--time and fluctuation rescalings, limiting random object, and mode of convergence. None of these is fixed by the supplied statement. In particular, it does not distinguish convergence of a single observable, convergence of finite-dimensional distributions, and convergence of a random field or process. It also does not say whether KPZ or EW behaviour is intended to mean scaling exponents, limiting distributions, or a particular stochastic evolution. Thus the supplied record states a research programme seeking sufficient hypotheses, rather than a fully quantified conjecture for a specified family of systems; a unique precise limit theorem cannot be extracted without additional source information.
\subsection{Short English statement}
Find conditions ensuring that integrated density fluctuations in one-dimensional interacting particle systems beyond nearest-neighbour exclusion, including finite-range exclusion, have KPZ or Edwards--Wilkinson scaling limits. The supplied record leaves the models, hypotheses, rescalings, and convergence targets unspecified, so a precise theorem remains to be formulated.
\subsection{Sources}
\begin{itemize}
\item Supplied local UnsolvedMath registry record, permanent ID \(30001618\), title: Universality Classes for Interacting Particle Fluctuations; exact source problem code: \texttt{catalogue entry}. Collection: Oberwolfach Reports Open Problems. No individual proposer is supplied; the proposed year is recorded as 2010.
\item Supplied report citation: \emph{Large Scale Stochastic Dynamics} (2011), Oberwolfach Reports, pp.~2877--2975, DOI \texttt{10.4171/owr/2010/50}. Supplied source URL and literature-source link, labelled ``original OWR report'': \url{https://doi.org/10.4171/owr/2010/50}. The background describes Oberwolfach Reports as documenting mathematical workshops and their open-problem discussions. No pinpoint page or verbatim report passage is supplied for this problem.
\item Literature assessment dated \texttt{2026-08-21}: the supplied curation labels the problem open, attributes to the report an explicit identification of rigorous KPZ and Edwards--Wilkinson universality beyond nearest-neighbour exclusion as an important unresolved problem, and states that later results cover special finite-range and integrable models rather than the general extension. This is supplied dated curation, not independently verified here. The record itself requests independent verification of time-sensitive status and cited claims; it supplies no separate references for those later results.
\item The supplied statement status is \texttt{corrected\_verified}; its accompanying verification note states only that TeX-style punctuation was normalized without changing the mathematical content. This metadata does not establish mathematical correctness or verify the literature assessment.
\end{itemize}
\end{document}
